%% kapa: exam=1 round=36 year=2025 period=1 form=A subject=civil_law subject_ko=민법 seq=34 number=34
%% source: Q-Net 감정평가사 기출문제 (cst00309) artlSeq=5231525 file=raw/1008/5231525/2025년 제36회 감정평가사 1차시험 1교시 문제지 원본.pdf
%% conversion: kordoc 4.24.51 -> tools/gen1.py (auto); review=proofread-sonnet; flags=none
\begin{question}{34}
창고소유자 甲은 乙과 월차임을 100만 원, 보증금을 3억 원으로 하는 창고임대차계약을 체결하고, 보증금반환청구권을 담보하기 위하여 乙에게 전세금을 3억 원, 임대차기간과 동일한 기간을 존속기간으로 하는 전세권등기도 경료해 주었다. 이에 관한 설명으로 옳은 것은? (다툼이 있으면 판례에 따름)
\begin{choices}
  \item 乙은 甲에게 전세권의 효력을 주장할 수 없고 임대차계약의 효력만 주장할 수 있다.
  \item 乙은 甲에게서 창고를 매수한 丙에게 자신의 전세권을 주장할 수 없다.
  \item 임대차계약 만료 후 甲이 보증금의 반환을 지체하면 乙은 전세권에 기한 경매를 신청할 수 없다.
  \item 임대차계약 만료 후 甲은 보증금에서 乙이 연체한 차임을 공제한 금액을 반환하면 된다.
  \item 乙이 丁에게 금전을 차용하면서 전세권저당권을 설정해 줄 당시 丁이 甲과 乙의 임대차 관계를 안 경우, 甲은 丁에게 乙의 연체차임으로 대항할 수 없다.
\end{choices}
\end{question}
\answer{4}
